\documentclass[11pt,a4paper]{article}
\usepackage{jheppub}
\title{Four-Loop Leading Singularities Summary}
\author{AntiGravity}
\abstract{This document summarizes the leading singularities of the 412 integrands. Integrands are classified by their set of leading singularities (and the number of such singularities). All failed or empty results are grouped into a single class.}
\begin{document}
\maketitle

\section{Summary of Results}
Total integrands: 412. Total distinct classes: 74.

\subsection*{Class 1 (185 integrands)}
Integrand Indices: 7, 8, 10, 11, 13, 14, 16, 26, 29, 33, 34, 42, 43, 51, 55, 57, 58, 63, 71, 73, 79, 83, 86, 87, 88, 90, 92, 99, 101, 102, 105, 109, 114, 115, 121, 124, 125, 127, 146, 148, 151, 152, 154, 156, 159, 161, 165, 166, 170, 185, 191, 196, 198, 199, 201, 205, 213, 216, 217, 222, 226, 227, 229, 231, 232, 234, 236, 240, 242, 247, 252, 253, 254, 257, 258, 259, 260, 261, 265, 267, 269, 273, 275, 276, 277, 278, 280, 282, 283, 284, 285, 286, 287, 288, 289, 291, 292, 293, 294, 296, 297, 301, 302, 303, 304, 305, 306, 307, 308, 309, 310, 311, 312, 313, 321, 324, 327, 328, 329, 330, 331, 334, 335, 337, 339, 341, 342, 343, 344, 345, 346, 347, 352, 353, 354, 355, 356, 357, 358, 359, 360, 361, 362, 364, 366, 367, 368, 369, 375, 376, 377, 378, 379, 380, 382, 383, 384, 385, 386, 387, 388, 389, 390, 391, 392, 393, 394, 395, 396, 397, 398, 399, 400, 401, 402, 403, 404, 405, 406, 407, 408, 409, 410, 411, 412

Result: Empty set $\{\}$ (Weight drop).\\[1em]

\subsection*{Class 2 (80 integrands)}
Integrand Indices: 1, 2, 4, 12, 17, 19, 20, 36, 40, 45, 47, 52, 54, 56, 61, 62, 65, 68, 76, 77, 82, 93, 97, 100, 103, 107, 129, 131, 136, 138, 143, 145, 149, 150, 153, 158, 172, 179, 180, 184, 186, 189, 192, 193, 194, 200, 204, 207, 210, 211, 215, 221, 224, 238, 248, 264, 268, 271, 281, 290, 298, 299, 314, 315, 316, 319, 320, 322, 323, 333, 338, 340, 351, 363, 370, 371, 372, 373, 374, 381

\textbf{Note:} Encountered uncancelled higher poles during evaluation.\\[1em]

Result: Empty set $\{\}$ (Weight drop).\\[1em]

\subsection*{Class 3 (41 integrands)}
Integrand Indices: 3, 5, 6, 9, 15, 18, 21, 22, 23, 32, 44, 60, 69, 70, 72, 74, 75, 78, 80, 81, 84, 85, 94, 95, 98, 104, 106, 108, 119, 120, 139, 160, 164, 174, 175, 195, 197, 206, 208, 235, 255

Number of Leading Singularities: 1

\begin{equation}
\frac{1}{\sqrt{u^2-2 u v-2 u+v^2-2 v+1}}
\end{equation}


\subsection*{Class 4 (11 integrands)}
Integrand Indices: 182, 202, 203, 209, 218, 219, 220, 225, 249, 317, 348

\textbf{Note:} Encountered uncancelled higher poles during evaluation.\\[1em]

Number of Leading Singularities: 1

\begin{equation}
\frac{1}{\sqrt{u^2-2 u v-2 u+v^2-2 v+1}}
\end{equation}


\subsection*{Class 5 (8 integrands)}
Integrand Indices: 59, 64, 96, 128, 134, 162, 181, 187

Number of Leading Singularities: 3

\begin{equation}
\frac{1}{\sqrt{u^2-2 u v-2 u+v^2-2 v+1}}
\end{equation}
\begin{equation}
\frac{1}{(v-1) \sqrt{u^2-2 u v-2 u+v^2-2 v+1}}
\end{equation}
\begin{equation}
\frac{v}{(v-1) \sqrt{u^2-2 u v-2 u+v^2-2 v+1}}
\end{equation}


\subsection*{Class 6 (6 integrands)}
Integrand Indices: 31, 35, 39, 46, 126, 133

Number of Leading Singularities: 1

\begin{equation}
\frac{v}{u^2-2 u v-2 u+v^2-2 v+1}
\end{equation}


\subsection*{Class 7 (4 integrands)}
Integrand Indices: 24, 27, 89, 122

Number of Leading Singularities: 1

\begin{equation}
\frac{u}{(v-1) \sqrt{u^2-2 u v-2 u+v^2-2 v+1}}
\end{equation}


\subsection*{Class 8 (4 integrands)}
Integrand Indices: 130, 142, 167, 246

Number of Leading Singularities: 2

\begin{equation}
\frac{v}{u^2-2 u v-2 u+v^2-2 v+1}
\end{equation}
\begin{equation}
\frac{v}{(v-1) \sqrt{u^2-2 u v-2 u+v^2-2 v+1}}
\end{equation}


\subsection*{Class 9 (4 integrands)}
Integrand Indices: 171, 177, 270, 332

\textbf{Note:} Encountered uncancelled higher poles during evaluation.\\[1em]

Number of Leading Singularities: 2

\begin{equation}
\frac{1}{\sqrt{u^2-2 u v-2 u+v^2-2 v+1}}
\end{equation}
\begin{equation}
\frac{u}{(v-1) \sqrt{u^2-2 u v-2 u+v^2-2 v+1}}
\end{equation}


\subsection*{Class 10 (3 integrands)}
Integrand Indices: 25, 28, 37

Number of Leading Singularities: 1

\begin{equation}
\frac{v}{(v-1) \sqrt{u^2-2 u v-2 u+v^2-2 v+1}}
\end{equation}


\subsection*{Class 11 (2 integrands)}
Integrand Indices: 30, 110

Number of Leading Singularities: 1

\begin{equation}
\frac{1}{v \sqrt{u^2-2 u v-2 u+v^2-2 v+1}}
\end{equation}


\subsection*{Class 12 (2 integrands)}
Integrand Indices: 38, 140

Number of Leading Singularities: 1

\begin{equation}
\frac{1}{u^2-2 u v-2 u+v^2-2 v+1}
\end{equation}


\subsection*{Class 13 (1 integrands)}
Integrand Indices: 41

Number of Leading Singularities: 1

\begin{equation}
\frac{v}{(u-v-1) \sqrt{u^2-2 u v-2 u+v^2-2 v+1}}
\end{equation}


\subsection*{Class 14 (1 integrands)}
Integrand Indices: 48

Number of Leading Singularities: 1

\begin{equation}
\frac{v}{(v-1) \left(u^2-2 u v-2 u+v^2-2 v+1\right)}
\end{equation}


\subsection*{Class 15 (1 integrands)}
Integrand Indices: 49

Number of Leading Singularities: 1

\begin{equation}
\frac{u}{v \left(u^2-2 u v-2 u+v^2-2 v+1\right)}
\end{equation}


\subsection*{Class 16 (1 integrands)}
Integrand Indices: 50

Number of Leading Singularities: 1

\begin{equation}
\frac{v \text{V2}(x(1,3)) \text{V2}(x(2,4))}{u^2-2 u v-2 u+v^2-2 v+1}
\end{equation}


\subsection*{Class 17 (1 integrands)}
Integrand Indices: 53

Number of Leading Singularities: 1

\begin{equation}
\frac{u v}{\left(u^2-2 u v-2 u+v^2-2 v+1\right)^2}
\end{equation}


\subsection*{Class 18 (1 integrands)}
Integrand Indices: 66

Number of Leading Singularities: 2

\begin{equation}
\frac{u-v-1}{u^2-2 u v-2 u+v^2-2 v+1}
\end{equation}
\begin{equation}
\frac{1}{\sqrt{u^2-2 u v-2 u+v^2-2 v+1}}
\end{equation}


\subsection*{Class 19 (1 integrands)}
Integrand Indices: 67

Number of Leading Singularities: 2

\begin{equation}
\frac{1}{u \sqrt{u^2-2 u v-2 u+v^2-2 v+1}}
\end{equation}
\begin{equation}
\frac{v}{u \sqrt{u^2-2 u v-2 u+v^2-2 v+1}}
\end{equation}


\subsection*{Class 20 (1 integrands)}
Integrand Indices: 91

Number of Leading Singularities: 3

\begin{equation}
\frac{u}{\sqrt{u^2-2 u v-2 u+v^2-2 v+1}}
\end{equation}
\begin{equation}
\frac{u}{(v-1) \sqrt{u^2-2 u v-2 u+v^2-2 v+1}}
\end{equation}
\begin{equation}
\frac{u v}{(v-1) \sqrt{u^2-2 u v-2 u+v^2-2 v+1}}
\end{equation}


\subsection*{Class 21 (1 integrands)}
Integrand Indices: 111

Number of Leading Singularities: 4

\begin{equation}
\frac{1}{u \sqrt{u^2-2 u v-2 u+v^2-2 v+1}}
\end{equation}
\begin{equation}
\frac{1}{(u-v) \sqrt{u^2-2 u v-2 u+v^2-2 v+1}}
\end{equation}
\begin{equation}
\frac{v}{u (u-v) \sqrt{u^2-2 u v-2 u+v^2-2 v+1}}
\end{equation}
\begin{equation}
\frac{1}{\sqrt{u^2-2 u v+v^2-4 v} \sqrt{u^2-2 u v-2 u+v^2-2 v+1}}
\end{equation}


\subsection*{Class 22 (1 integrands)}
Integrand Indices: 112

Number of Leading Singularities: 2

\begin{equation}
\frac{1}{(u-v) \sqrt{u^2-2 u v-2 u+v^2-2 v+1}}
\end{equation}
\begin{equation}
\frac{1}{v \sqrt{u^2-2 u v-2 u+v^2-2 v+1}}
\end{equation}


\subsection*{Class 23 (1 integrands)}
Integrand Indices: 113

Number of Leading Singularities: 2

\begin{equation}
\frac{v (u-v+1)}{u \left(u^2-2 u v-2 u+v^2-2 v+1\right)}
\end{equation}
\begin{equation}
\frac{v}{u \sqrt{u^2-2 u v-2 u+v^2-2 v+1}}
\end{equation}


\subsection*{Class 24 (1 integrands)}
Integrand Indices: 116

Number of Leading Singularities: 2

\begin{equation}
\frac{1}{u^2-2 u v-2 u+v^2-2 v+1}
\end{equation}
\begin{equation}
\frac{1}{(u-1) \sqrt{u^2-2 u v-2 u+v^2-2 v+1}}
\end{equation}


\subsection*{Class 25 (1 integrands)}
Integrand Indices: 117

Number of Leading Singularities: 1

\begin{equation}
\frac{1}{(u-1) \sqrt{u^2-2 u v-2 u+v^2-2 v+1}}
\end{equation}


\subsection*{Class 26 (1 integrands)}
Integrand Indices: 118

Number of Leading Singularities: 5

\begin{equation}
\frac{1}{(u-1) \sqrt{u^2-2 u v-2 u+v^2-2 v+1}}
\end{equation}
\begin{equation}
\frac{1}{(u-v-1) \sqrt{u^2-2 u v-2 u+v^2-2 v+1}}
\end{equation}
\begin{equation}
\frac{v}{(u-1) (u-v-1) \sqrt{u^2-2 u v-2 u+v^2-2 v+1}}
\end{equation}
\begin{equation}
\frac{1}{(u+v-1) \sqrt{u^2-2 u v-2 u+v^2-2 v+1}}
\end{equation}
\begin{equation}
\frac{v}{(u-1) (u+v-1) \sqrt{u^2-2 u v-2 u+v^2-2 v+1}}
\end{equation}


\subsection*{Class 27 (1 integrands)}
Integrand Indices: 123

Number of Leading Singularities: 2

\begin{equation}
\frac{v}{u^2-2 u v-2 u+v^2-2 v+1}
\end{equation}
\begin{equation}
\frac{v}{(u-v) \sqrt{u^2-2 u v-2 u+v^2-2 v+1}}
\end{equation}


\subsection*{Class 28 (1 integrands)}
Integrand Indices: 132

Number of Leading Singularities: 3

\begin{equation}
\frac{v}{u^2-2 u v-2 u+v^2-2 v+1}
\end{equation}
\begin{equation}
\frac{1}{\sqrt{u^2-2 u v-2 u+v^2-2 v+1}}
\end{equation}
\begin{equation}
\frac{v}{(v-1) \sqrt{u^2-2 u v-2 u+v^2-2 v+1}}
\end{equation}


\subsection*{Class 29 (1 integrands)}
Integrand Indices: 135

Number of Leading Singularities: 5

\begin{equation}
\frac{1}{u^2-2 u v-2 u+v^2-2 v+1}
\end{equation}
\begin{equation}
\frac{v}{u^2-2 u v-2 u+v^2-2 v+1}
\end{equation}
\begin{equation}
\frac{1}{\sqrt{u^2-2 u v-2 u+v^2-2 v+1}}
\end{equation}
\begin{equation}
\frac{1}{(v-1) \sqrt{u^2-2 u v-2 u+v^2-2 v+1}}
\end{equation}
\begin{equation}
\frac{v}{(v-1) \sqrt{u^2-2 u v-2 u+v^2-2 v+1}}
\end{equation}


\subsection*{Class 30 (1 integrands)}
Integrand Indices: 137

Number of Leading Singularities: 2

\begin{equation}
\frac{u}{(v-1)^2 \sqrt{u^2-2 u v-2 u+v^2-2 v+1}}
\end{equation}
\begin{equation}
\frac{u v}{(v-1)^2 \sqrt{u^2-2 u v-2 u+v^2-2 v+1}}
\end{equation}


\subsection*{Class 31 (1 integrands)}
Integrand Indices: 141

Number of Leading Singularities: 3

\begin{equation}
\frac{u}{u^2-2 u v-2 u+v^2-2 v+1}
\end{equation}
\begin{equation}
\frac{u}{(v-1) \left(u^2-2 u v-2 u+v^2-2 v+1\right)}
\end{equation}
\begin{equation}
\frac{u v}{(v-1) \left(u^2-2 u v-2 u+v^2-2 v+1\right)}
\end{equation}


\subsection*{Class 32 (1 integrands)}
Integrand Indices: 144

Number of Leading Singularities: 3

\begin{equation}
\frac{1}{u^2-2 u v-2 u+v^2-2 v+1}
\end{equation}
\begin{equation}
\frac{1}{(u-1) \sqrt{u^2-2 u v-2 u+v^2-2 v+1}}
\end{equation}
\begin{equation}
\frac{1}{(u-v-1) \sqrt{u^2-2 u v-2 u+v^2-2 v+1}}
\end{equation}


\subsection*{Class 33 (1 integrands)}
Integrand Indices: 147

Number of Leading Singularities: 4

\begin{equation}
\frac{u}{v \left(u^2-2 u v-2 u+v^2-2 v+1\right)}
\end{equation}
\begin{equation}
\frac{u}{(u-1) v \sqrt{u^2-2 u v-2 u+v^2-2 v+1}}
\end{equation}
\begin{equation}
\frac{u}{v (u-v-1) \sqrt{u^2-2 u v-2 u+v^2-2 v+1}}
\end{equation}
\begin{equation}
\frac{u}{v (u+v-1) \sqrt{u^2-2 u v-2 u+v^2-2 v+1}}
\end{equation}


\subsection*{Class 34 (1 integrands)}
Integrand Indices: 155

Number of Leading Singularities: 4

\begin{equation}
\frac{v}{u^2-2 u v-2 u+v^2-2 v+1}
\end{equation}
\begin{equation}
\frac{v}{(u-v) \sqrt{u^2-2 u v-2 u+v^2-2 v+1}}
\end{equation}
\begin{equation}
\frac{v}{(u-v+1) \sqrt{u^2-2 u v-2 u+v^2-2 v+1}}
\end{equation}
\begin{equation}
\frac{v}{(v-1) \sqrt{u^2-2 u v-2 u+v^2-2 v+1}}
\end{equation}


\subsection*{Class 35 (1 integrands)}
Integrand Indices: 157

Number of Leading Singularities: 2

\begin{equation}
\frac{v}{\left(u^2-2 u v-2 u+v^2-2 v+1\right)^{3/2}}
\end{equation}
\begin{equation}
\frac{v}{(v-1) \left(u^2-2 u v-2 u+v^2-2 v+1\right)}
\end{equation}


\subsection*{Class 36 (1 integrands)}
Integrand Indices: 163

Number of Leading Singularities: 10

\begin{equation}
\frac{1}{u^2-2 u v-2 u+v^2-2 v+1}
\end{equation}
\begin{equation}
\frac{u}{u^2-2 u v-2 u+v^2-2 v+1}
\end{equation}
\begin{equation}
\frac{v}{u^2-2 u v-2 u+v^2-2 v+1}
\end{equation}
\begin{equation}
\frac{1}{\sqrt{u^2-2 u v-2 u+v^2-2 v+1}}
\end{equation}
\begin{equation}
\frac{1}{(u-1) \sqrt{u^2-2 u v-2 u+v^2-2 v+1}}
\end{equation}
\begin{equation}
\frac{u}{(u-1) \sqrt{u^2-2 u v-2 u+v^2-2 v+1}}
\end{equation}
\begin{equation}
\frac{u}{(u-v) \sqrt{u^2-2 u v-2 u+v^2-2 v+1}}
\end{equation}
\begin{equation}
\frac{1}{(v-1) \sqrt{u^2-2 u v-2 u+v^2-2 v+1}}
\end{equation}
\begin{equation}
\frac{v}{(u-v) \sqrt{u^2-2 u v-2 u+v^2-2 v+1}}
\end{equation}
\begin{equation}
\frac{v}{(v-1) \sqrt{u^2-2 u v-2 u+v^2-2 v+1}}
\end{equation}


\subsection*{Class 37 (1 integrands)}
Integrand Indices: 168

Number of Leading Singularities: 4

\begin{equation}
\frac{v}{u^2-2 u v-2 u+v^2-2 v+1}
\end{equation}
\begin{equation}
\frac{1}{\sqrt{u^2-2 u v-2 u+v^2-2 v+1}}
\end{equation}
\begin{equation}
\frac{v}{(u-v) \sqrt{u^2-2 u v-2 u+v^2-2 v+1}}
\end{equation}
\begin{equation}
\frac{v}{(v-1) \sqrt{u^2-2 u v-2 u+v^2-2 v+1}}
\end{equation}


\subsection*{Class 38 (1 integrands)}
Integrand Indices: 169

Number of Leading Singularities: 2

\begin{equation}
\frac{v}{u \left(u^2-2 u v-2 u+v^2-2 v+1\right)}
\end{equation}
\begin{equation}
\frac{v}{u (v-1) \sqrt{u^2-2 u v-2 u+v^2-2 v+1}}
\end{equation}


\subsection*{Class 39 (1 integrands)}
Integrand Indices: 173

Number of Leading Singularities: 2

\begin{equation}
\frac{u-v+1}{u^2-2 u v-2 u+v^2-2 v+1}
\end{equation}
\begin{equation}
\frac{1}{\sqrt{u^2-2 u v-2 u+v^2-2 v+1}}
\end{equation}


\subsection*{Class 40 (1 integrands)}
Integrand Indices: 176

Number of Leading Singularities: 2

\begin{equation}
\frac{1}{\sqrt{u^2-2 u v-2 u+v^2-2 v+1}}
\end{equation}
\begin{equation}
\frac{u}{(u-1) \sqrt{u^2-2 u v-2 u+v^2-2 v+1}}
\end{equation}


\subsection*{Class 41 (1 integrands)}
Integrand Indices: 178

Number of Leading Singularities: 2

\begin{equation}
\frac{1}{\sqrt{u^2-2 u v-2 u+v^2-2 v+1}}
\end{equation}
\begin{equation}
\frac{1}{(u-v) \sqrt{u^2-2 u v-2 u+v^2-2 v+1}}
\end{equation}


\subsection*{Class 42 (1 integrands)}
Integrand Indices: 183

\textbf{Note:} Encountered uncancelled higher poles during evaluation.\\[1em]

Number of Leading Singularities: 2

\begin{equation}
\frac{v}{u^2-2 u v-2 u+v^2-2 v+1}
\end{equation}
\begin{equation}
\frac{1}{\sqrt{u^2-2 u v-2 u+v^2-2 v+1}}
\end{equation}


\subsection*{Class 43 (1 integrands)}
Integrand Indices: 188

Number of Leading Singularities: 2

\begin{equation}
\frac{1}{\sqrt{u^2-2 u v-2 u+v^2-2 v+1}}
\end{equation}
\begin{equation}
\frac{v}{\sqrt{u^2-2 u v-2 u+v^2-2 v+1}}
\end{equation}


\subsection*{Class 44 (1 integrands)}
Integrand Indices: 190

Number of Leading Singularities: 3

\begin{equation}
\frac{1}{\sqrt{u^2-2 u v-2 u+v^2-2 v+1}}
\end{equation}
\begin{equation}
\frac{1}{(u-1) \sqrt{u^2-2 u v-2 u+v^2-2 v+1}}
\end{equation}
\begin{equation}
\frac{1}{v \sqrt{u^2-2 u v-2 u+v^2-2 v+1}}
\end{equation}


\subsection*{Class 45 (1 integrands)}
Integrand Indices: 212

Number of Leading Singularities: 4

\begin{equation}
\frac{u-v-1}{u^2-2 u v-2 u+v^2-2 v+1}
\end{equation}
\begin{equation}
\frac{1}{\sqrt{u^2-2 u v-2 u+v^2-2 v+1}}
\end{equation}
\begin{equation}
\frac{1}{(v-1) \sqrt{u^2-2 u v-2 u+v^2-2 v+1}}
\end{equation}
\begin{equation}
\frac{v}{(v-1) \sqrt{u^2-2 u v-2 u+v^2-2 v+1}}
\end{equation}


\subsection*{Class 46 (1 integrands)}
Integrand Indices: 214

Number of Leading Singularities: 6

\begin{equation}
\frac{v}{\sqrt{u^2-2 u v-2 u+v^2-2 v+1}}
\end{equation}
\begin{equation}
\frac{v}{(u-1) \sqrt{u^2-2 u v-2 u+v^2-2 v+1}}
\end{equation}
\begin{equation}
\frac{u v}{(u-1) \sqrt{u^2-2 u v-2 u+v^2-2 v+1}}
\end{equation}
\begin{equation}
\frac{u v}{\sqrt{u^2-2 u-4 v+1} \sqrt{u^2-2 u v-2 u+v^2-2 v+1}}
\end{equation}
\begin{equation}
\frac{v}{(v-1) \sqrt{u^2-2 u v-2 u+v^2-2 v+1}}
\end{equation}
\begin{equation}
\frac{v^2}{(v-1) \sqrt{u^2-2 u v-2 u+v^2-2 v+1}}
\end{equation}


\subsection*{Class 47 (1 integrands)}
Integrand Indices: 223

\textbf{Note:} Encountered uncancelled higher poles during evaluation.\\[1em]

Number of Leading Singularities: 2

\begin{equation}
\frac{1}{\sqrt{u^2-2 u v-2 u+v^2-2 v+1}}
\end{equation}
\begin{equation}
\frac{v}{(u-v-1) \sqrt{u^2-2 u v-2 u+v^2-2 v+1}}
\end{equation}


\subsection*{Class 48 (1 integrands)}
Integrand Indices: 228

Number of Leading Singularities: 2

\begin{equation}
\frac{u}{(u-1) \sqrt{u^2-2 u v-2 u+v^2-2 v+1}}
\end{equation}
\begin{equation}
\frac{u}{v \sqrt{u^2-2 u v-2 u+v^2-2 v+1}}
\end{equation}


\subsection*{Class 49 (1 integrands)}
Integrand Indices: 230

\textbf{Note:} Encountered uncancelled higher poles during evaluation.\\[1em]

Number of Leading Singularities: 1

\begin{equation}
\frac{1}{u^2-2 u v-2 u+v^2-2 v+1}
\end{equation}


\subsection*{Class 50 (1 integrands)}
Integrand Indices: 233

Number of Leading Singularities: 4

\begin{equation}
\frac{u}{u^2-2 u v-2 u+v^2-2 v+1}
\end{equation}
\begin{equation}
\frac{u}{(u-1) \sqrt{u^2-2 u v-2 u+v^2-2 v+1}}
\end{equation}
\begin{equation}
\frac{u}{(v-1) \sqrt{u^2-2 u v-2 u+v^2-2 v+1}}
\end{equation}
\begin{equation}
\frac{u}{(u+v-1) \sqrt{u^2-2 u v-2 u+v^2-2 v+1}}
\end{equation}


\subsection*{Class 51 (1 integrands)}
Integrand Indices: 237

\textbf{Note:} Encountered uncancelled higher poles during evaluation.\\[1em]

Number of Leading Singularities: 2

\begin{equation}
\frac{1}{\sqrt{u^2-2 u v-2 u+v^2-2 v+1}}
\end{equation}
\begin{equation}
\frac{u}{(u-1) \sqrt{u^2-2 u v-2 u+v^2-2 v+1}}
\end{equation}


\subsection*{Class 52 (1 integrands)}
Integrand Indices: 239

\textbf{Note:} Encountered uncancelled higher poles during evaluation.\\[1em]

Number of Leading Singularities: 2

\begin{equation}
\frac{v}{u \sqrt{u^2-2 u v-2 u+v^2-2 v+1}}
\end{equation}
\begin{equation}
\frac{v}{(v-1) \sqrt{u^2-2 u v-2 u+v^2-2 v+1}}
\end{equation}


\subsection*{Class 53 (1 integrands)}
Integrand Indices: 241

Number of Leading Singularities: 3

\begin{equation}
\frac{1}{\sqrt{u^2-2 u v-2 u+v^2-2 v+1}}
\end{equation}
\begin{equation}
\frac{u-v}{u \sqrt{u^2-2 u v-2 u+v^2-2 v+1}}
\end{equation}
\begin{equation}
\frac{v}{u \sqrt{u^2-2 u v-2 u+v^2-2 v+1}}
\end{equation}


\subsection*{Class 54 (1 integrands)}
Integrand Indices: 243

Number of Leading Singularities: 3

\begin{equation}
\frac{v}{(u-1) \sqrt{u^2-2 u v-2 u+v^2-2 v+1}}
\end{equation}
\begin{equation}
\frac{v}{(u-v) \sqrt{u^2-2 u v-2 u+v^2-2 v+1}}
\end{equation}
\begin{equation}
\frac{v}{(v-1) \sqrt{u^2-2 u v-2 u+v^2-2 v+1}}
\end{equation}


\subsection*{Class 55 (1 integrands)}
Integrand Indices: 244

Number of Leading Singularities: 3

\begin{equation}
\frac{v}{u^2-2 u v-2 u+v^2-2 v+1}
\end{equation}
\begin{equation}
\frac{1}{\sqrt{u^2-2 u v-2 u+v^2-2 v+1}}
\end{equation}
\begin{equation}
\frac{v}{\sqrt{u^2-2 u v-2 u+v^2-2 v+1}}
\end{equation}


\subsection*{Class 56 (1 integrands)}
Integrand Indices: 245

Number of Leading Singularities: 2

\begin{equation}
\frac{1}{\sqrt{u^2-2 u v-2 u+v^2-2 v+1}}
\end{equation}
\begin{equation}
\frac{v}{(u-1) \sqrt{u^2-2 u v-2 u+v^2-2 v+1}}
\end{equation}


\subsection*{Class 57 (1 integrands)}
Integrand Indices: 250

Number of Leading Singularities: 3

\begin{equation}
\frac{1}{\sqrt{u^2-2 u v-2 u+v^2-2 v+1}}
\end{equation}
\begin{equation}
\frac{u}{(v-1) \sqrt{u^2-2 u v-2 u+v^2-2 v+1}}
\end{equation}
\begin{equation}
\frac{u}{v \sqrt{u^2-2 u v-2 u+v^2-2 v+1}}
\end{equation}


\subsection*{Class 58 (1 integrands)}
Integrand Indices: 251

Number of Leading Singularities: 2

\begin{equation}
\frac{1}{u^2-2 u v-2 u+v^2-2 v+1}
\end{equation}
\begin{equation}
\frac{v}{u^2-2 u v-2 u+v^2-2 v+1}
\end{equation}


\subsection*{Class 59 (1 integrands)}
Integrand Indices: 256

\textbf{Note:} Encountered uncancelled higher poles during evaluation.\\[1em]

Number of Leading Singularities: 2

\begin{equation}
\frac{v}{u^2-2 u v-2 u+v^2-2 v+1}
\end{equation}
\begin{equation}
\frac{u v}{(v-1) \left(u^2-2 u v-2 u+v^2-2 v+1\right)}
\end{equation}


\subsection*{Class 60 (1 integrands)}
Integrand Indices: 262

Number of Leading Singularities: 4

\begin{equation}
\frac{u}{u^2-2 u v-2 u+v^2-2 v+1}
\end{equation}
\begin{equation}
\frac{u (u-v-1)}{(u-v+1) \left(u^2-2 u v-2 u+v^2-2 v+1\right)}
\end{equation}
\begin{equation}
\frac{1}{\sqrt{u^2-2 u v-2 u+v^2-2 v+1}}
\end{equation}
\begin{equation}
\frac{u}{(v-1) \sqrt{u^2-2 u v-2 u+v^2-2 v+1}}
\end{equation}


\subsection*{Class 61 (1 integrands)}
Integrand Indices: 263

\textbf{Note:} Encountered uncancelled higher poles during evaluation.\\[1em]

Number of Leading Singularities: 2

\begin{equation}
\frac{u}{u^2-2 u v-2 u+v^2-2 v+1}
\end{equation}
\begin{equation}
\frac{1}{\sqrt{u^2-2 u v-2 u+v^2-2 v+1}}
\end{equation}


\subsection*{Class 62 (1 integrands)}
Integrand Indices: 266

\textbf{Note:} Encountered uncancelled higher poles during evaluation.\\[1em]

Number of Leading Singularities: 8

\begin{equation}
\frac{1}{\sqrt{u^2-2 u v-2 u+v^2-2 v+1}}
\end{equation}
\begin{equation}
\frac{u}{(u-1) \sqrt{u^2-2 u v-2 u+v^2-2 v+1}}
\end{equation}
\begin{equation}
\frac{u}{(u-v-1) \sqrt{u^2-2 u v-2 u+v^2-2 v+1}}
\end{equation}
\begin{equation}
\frac{u}{(u-v) \sqrt{u^2-2 u v-2 u+v^2-2 v+1}}
\end{equation}
\begin{equation}
\frac{u v}{(u-1) (u-v-1) \sqrt{u^2-2 u v-2 u+v^2-2 v+1}}
\end{equation}
\begin{equation}
\frac{v}{(u-v) \sqrt{u^2-2 u v-2 u+v^2-2 v+1}}
\end{equation}
\begin{equation}
\frac{u}{(u+v-1) \sqrt{u^2-2 u v-2 u+v^2-2 v+1}}
\end{equation}
\begin{equation}
\frac{u v}{(u-1) (u+v-1) \sqrt{u^2-2 u v-2 u+v^2-2 v+1}}
\end{equation}


\subsection*{Class 63 (1 integrands)}
Integrand Indices: 272

\textbf{Note:} Encountered uncancelled higher poles during evaluation.\\[1em]

Number of Leading Singularities: 2

\begin{equation}
\frac{1}{\sqrt{u^2-2 u v-2 u+v^2-2 v+1}}
\end{equation}
\begin{equation}
\frac{u}{v \sqrt{u^2-2 u v-2 u+v^2-2 v+1}}
\end{equation}


\subsection*{Class 64 (1 integrands)}
Integrand Indices: 274

\textbf{Note:} Encountered uncancelled higher poles during evaluation.\\[1em]

Number of Leading Singularities: 4

\begin{equation}
\frac{1}{\sqrt{u^2-2 u v-2 u+v^2-2 v+1}}
\end{equation}
\begin{equation}
\frac{u}{(u-1) \sqrt{u^2-2 u v-2 u+v^2-2 v+1}}
\end{equation}
\begin{equation}
\frac{u}{(v-1) \sqrt{u^2-2 u v-2 u+v^2-2 v+1}}
\end{equation}
\begin{equation}
\frac{u}{v \sqrt{u^2-2 u v-2 u+v^2-2 v+1}}
\end{equation}


\subsection*{Class 65 (1 integrands)}
Integrand Indices: 279

\textbf{Note:} Encountered uncancelled higher poles during evaluation.\\[1em]

Number of Leading Singularities: 3

\begin{equation}
\frac{1}{\sqrt{u^2-2 u v-2 u+v^2-2 v+1}}
\end{equation}
\begin{equation}
\frac{1}{(v-1) \sqrt{u^2-2 u v-2 u+v^2-2 v+1}}
\end{equation}
\begin{equation}
\frac{v}{(v-1) \sqrt{u^2-2 u v-2 u+v^2-2 v+1}}
\end{equation}


\subsection*{Class 66 (1 integrands)}
Integrand Indices: 295

\textbf{Note:} Encountered uncancelled higher poles during evaluation.\\[1em]

Number of Leading Singularities: 2

\begin{equation}
\frac{1}{\sqrt{u^2-2 u v-2 u+v^2-2 v+1}}
\end{equation}
\begin{equation}
\frac{v}{(u-1) \sqrt{u^2-2 u v-2 u+v^2-2 v+1}}
\end{equation}


\subsection*{Class 67 (1 integrands)}
Integrand Indices: 300

\textbf{Note:} Encountered uncancelled higher poles during evaluation.\\[1em]

Number of Leading Singularities: 3

\begin{equation}
\frac{u}{u^2-2 u v-2 u+v^2-2 v+1}
\end{equation}
\begin{equation}
\frac{1}{\sqrt{u^2-2 u v-2 u+v^2-2 v+1}}
\end{equation}
\begin{equation}
\frac{u}{(v-1) \sqrt{u^2-2 u v-2 u+v^2-2 v+1}}
\end{equation}


\subsection*{Class 68 (1 integrands)}
Integrand Indices: 318

\textbf{Note:} Encountered uncancelled higher poles during evaluation.\\[1em]

Number of Leading Singularities: 2

\begin{equation}
\frac{1}{\sqrt{u^2-2 u v-2 u+v^2-2 v+1}}
\end{equation}
\begin{equation}
\frac{v}{(u+v-1) \sqrt{u^2-2 u v-2 u+v^2-2 v+1}}
\end{equation}


\subsection*{Class 69 (1 integrands)}
Integrand Indices: 325

Number of Leading Singularities: 5

\begin{equation}
\frac{1}{\sqrt{u^2-2 u v-2 u+v^2-2 v+1}}
\end{equation}
\begin{equation}
\frac{u}{(v-1)^2 \sqrt{u^2-2 u v-2 u+v^2-2 v+1}}
\end{equation}
\begin{equation}
\frac{u}{v \sqrt{u^2-2 u v-2 u+v^2-2 v+1}}
\end{equation}
\begin{equation}
\frac{u}{(v-1)^2 v \sqrt{u^2-2 u v-2 u+v^2-2 v+1}}
\end{equation}
\begin{equation}
\frac{v-u}{v \sqrt{u^2-2 u v-2 u+v^2-2 v+1}}
\end{equation}


\subsection*{Class 70 (1 integrands)}
Integrand Indices: 326

\textbf{Note:} Encountered uncancelled higher poles during evaluation.\\[1em]

Number of Leading Singularities: 7

\begin{equation}
\frac{1}{\sqrt{u^2-2 u v-2 u+v^2-2 v+1}}
\end{equation}
\begin{equation}
\frac{1}{(u-1) \sqrt{u^2-2 u v-2 u+v^2-2 v+1}}
\end{equation}
\begin{equation}
\frac{1}{u \sqrt{u^2-2 u v-2 u+v^2-2 v+1}}
\end{equation}
\begin{equation}
\frac{1}{(u-1) u \sqrt{u^2-2 u v-2 u+v^2-2 v+1}}
\end{equation}
\begin{equation}
\frac{1}{\sqrt{u^2-2 u-4 v+1} \sqrt{u^2-2 u v-2 u+v^2-2 v+1}}
\end{equation}
\begin{equation}
\frac{u-v}{u \sqrt{u^2-2 u v-2 u+v^2-2 v+1}}
\end{equation}
\begin{equation}
\frac{v}{u \sqrt{u^2-2 u v-2 u+v^2-2 v+1}}
\end{equation}


\subsection*{Class 71 (1 integrands)}
Integrand Indices: 336

\textbf{Note:} Encountered uncancelled higher poles during evaluation.\\[1em]

Number of Leading Singularities: 3

\begin{equation}
\frac{u-v+1}{u^2-2 u v-2 u+v^2-2 v+1}
\end{equation}
\begin{equation}
\frac{1}{\sqrt{u^2-2 u v-2 u+v^2-2 v+1}}
\end{equation}
\begin{equation}
\frac{1}{(u-v) \sqrt{u^2-2 u v-2 u+v^2-2 v+1}}
\end{equation}


\subsection*{Class 72 (1 integrands)}
Integrand Indices: 349

\textbf{Note:} Encountered uncancelled higher poles during evaluation.\\[1em]

Number of Leading Singularities: 4

\begin{equation}
\frac{v}{(u-1) \sqrt{u^2-2 u v-2 u+v^2-2 v+1}}
\end{equation}
\begin{equation}
\frac{v}{u \sqrt{u^2-2 u v-2 u+v^2-2 v+1}}
\end{equation}
\begin{equation}
\frac{v}{(u-v) \sqrt{u^2-2 u v-2 u+v^2-2 v+1}}
\end{equation}
\begin{equation}
\frac{v}{(v-1) \sqrt{u^2-2 u v-2 u+v^2-2 v+1}}
\end{equation}


\subsection*{Class 73 (1 integrands)}
Integrand Indices: 350

\textbf{Note:} Encountered uncancelled higher poles during evaluation.\\[1em]

Number of Leading Singularities: 2

\begin{equation}
\frac{1}{u^2-2 u v-2 u+v^2-2 v+1}
\end{equation}
\begin{equation}
\frac{1}{(u-1) \sqrt{u^2-2 u v-2 u+v^2-2 v+1}}
\end{equation}


\subsection*{Class 74 (1 integrands)}
Integrand Indices: 365

\textbf{Note:} Encountered uncancelled higher poles during evaluation.\\[1em]

Number of Leading Singularities: 4

\begin{equation}
\frac{1}{\sqrt{u^2-2 u v-2 u+v^2-2 v+1}}
\end{equation}
\begin{equation}
\frac{v}{(u-1) \sqrt{u^2-2 u v-2 u+v^2-2 v+1}}
\end{equation}
\begin{equation}
\frac{v}{(u-v-1) \sqrt{u^2-2 u v-2 u+v^2-2 v+1}}
\end{equation}
\begin{equation}
\frac{v}{(u+v-1) \sqrt{u^2-2 u v-2 u+v^2-2 v+1}}
\end{equation}


\end{document}
